\documentclass{book}
\usepackage[backend=biber,natbib=true,style=authoryear]{biblatex}
\addbibresource{/home/nguyen/1_NQBH/math/bib.bib}
\usepackage[utf8]{inputenc}
\usepackage{float}
\usepackage{graphicx}
\usepackage[colorlinks=true,linkcolor=blue,urlcolor=red,citecolor=magenta]{hyperref}
\usepackage{amsmath,amssymb,amsthm}
\allowdisplaybreaks
\numberwithin{equation}{section}
\newtheorem{assumption}{Assumption}[section]
\newtheorem{lemma}{Lemma}[section]
\newtheorem{corollary}{Corollary}[section]
\newtheorem{definition}{Definition}[section]
\newtheorem{proposition}{Proposition}[section]
\newtheorem{theorem}{Theorem}[section]
\newtheorem{notation}{Notation}[section]
\newtheorem{remark}{Remark}[section]
\newtheorem{example}{Example}[section]
\newtheorem{ques}{Question}[section]
\newtheorem{problem}{Problem}[section]
\newtheorem{conjecture}{Conjecture}[section]
\usepackage[left=0.5in,right=0.5in,top=1.5cm,bottom=1.5cm]{geometry}
\usepackage{fancyhdr}
\pagestyle{fancy}
\fancyhf{}
\addtolength{\headheight}{0pt}% obsolete
\lhead{\scriptsize \chaptername~\thechapter}
\rhead{\scriptsize \nouppercase{\leftmark}} %\nouppercase !
\renewcommand{\chaptermark}[1]{\markboth{#1}{}}
\cfoot{\thepage}
\def\labelitemii{$\circ$}

\title{Collins COBUILD English Grammar, 4the edition (2017) }
\author{Collector: Nguyen Quan Ba Hong}
\date{\today}

\begin{document}
\maketitle
\setcounter{secnumdepth}{6}
\setcounter{tocdepth}{6}
\tableofcontents

%------------------------------------------------------------------------------%

\section*{About COBUILD}
When the 1st COBUILD dictionary was published in 1987, it revolutionized dictionaries for learners.

It was the 1st of the new generation of language reference materials that were based on actual evidence of how English was used, rather than lexicographer intuition.

%
Collins and the University of Birmingham, led by the linguist John Sinclair, developed an electronic corpus in 1980s, called the Collins Birmingham University International Language Database (COBUILD).

This corpus, which for several years was known as the Bank of English\textcopyright, became the largest collection of English data in the world.

COBUILD dictionary editors use the corpus to analyze the way that people really use the language.

%
The Collins Corpus now contains 4.5 billion words taken from websites, newspapers, magazines and books published around the world, and from spoken material from radio, TV and everyday conversations.

New data is added to the corpus every month, to help COBUILD editors identify new words, grammatical structures, and meanings from the moment they are 1st used.

%
All COBUILD language reference books are based on the information their editors find the Collins Corpus.

Because the corpus is so large, their editors can look at lots of examples of how people really use the language.

The data tells them how the language is used; the function of different structures; which words are used together; and how often these words and structures are used.

%
All of the examples in COBUILD language materials are examples of real English, taken from the corpus.

The examples have been carefully chosen to demonstrate typical grammatical patterns, typically vocabulary and typical contexts.

%
COBUILD English Grammar is no exception: Collins editors and researchers have been able to use this wealthy of information to establish a unique and full description of English grammar, and to track the development of certain grammatical structures over time.

%
The corpus lies at the heart of COBUILD, and you can be confident that COBUILD will show you what you need to know to be able to communicate easily and accurately in English.

%
To learn more about COBUILD and the Collins Corpus, go to \url{www.collinselt.com} and click on `COBUILD Reference'.

%------------------------------------------------------------------------------%

\section*{Introduction}

\section*{How to use this Grammar}

\section*{Glossary of grammatical terms}

%------------------------------------------------------------------------------%

\chapter{Referring to people \& things}

%------------------------------------------------------------------------------%

\chapter{Giving information about people \& things}

%------------------------------------------------------------------------------%

\chapter{Types of verb}

%------------------------------------------------------------------------------%

\chapter{Expressing time: tenses \& time adverbials}

%------------------------------------------------------------------------------%

\chapter{Modals, negatives, \& ways of forming sentences}

%------------------------------------------------------------------------------%

\chapter{Expressing manner \& place}

%------------------------------------------------------------------------------%

\chapter{Reporting what people say or think}

%------------------------------------------------------------------------------%

\chapter{Combining messages}

%------------------------------------------------------------------------------%

\chapter{Changing the focus in a sentenec}

%------------------------------------------------------------------------------%

\chapter{Making a text hold together}

%------------------------------------------------------------------------------%

\chapter{Reference section}

%------------------------------------------------------------------------------%

\chapter{The grammar of business English}

%------------------------------------------------------------------------------%

\chapter{The grammar of academic English}

%------------------------------------------------------------------------------%

\chapter{About the Publisher}

%------------------------------------------------------------------------------%


\printbibliography[heading=bibintoc]
\end{document}